\section{Математические модели в биологии и биоинформатике}
\label{sec:models-in-biology-and-bioinformatics}

Математические модели в биологии описывают изменение численности
популяций, взаимодействие видов, эпидемические процессы, реакционно-
диффузионные системы и другие динамические явления. В биоинформатике
математические методы применяются к последовательностям ДНК, РНК и белков,
к филогенетическим данным и большим биологическим наборам данных.

\subsection{Модели роста популяции}

Простейшая модель Мальтуса предполагает постоянную удельную скорость роста:
\[
\boxed{\frac{dN}{dt}=rN},
\qquad
N(t)=N_0e^{rt}.
\]
Она пригодна лишь пока ограниченность ресурсов несущественна.

Модель Ферхюльста учитывает насыщение:
\[
\boxed{
\frac{dN}{dt}=rN\left(1-\frac{N}{K}\right),
}
\]
где $K$ --- ёмкость среды. Стационарные состояния:
\[
N_1^*=0,
\qquad
N_2^*=K.
\]
При $r>0$ состояние $N=0$ неустойчиво, а $N=K$ асимптотически устойчиво.
Таким образом, даже простая нелинейность качественно меняет долгосрочный
прогноз.

\subsection{Взаимодействие видов: модель Лотки--Вольтерры}

Классическая модель ``хищник--жертва'' имеет вид
\[
\boxed{
\begin{aligned}
\dot x&=ax-bxy,\\
\dot y&=-cy+dxy,
\end{aligned}
}
\qquad a,b,c,d>0,
\]
где $x$ --- численность жертв, $y$ --- хищников.

Кроме тривиального равновесия существует положительное стационарное
состояние
\[
\boxed{
(x^*,y^*)=\left(\frac{c}{d},\frac{a}{b}\right).
}
\]
Линеаризация системы около стационарного состояния позволяет исследовать
его устойчивость. В идеальной классической модели возникают нейтральные
колебания; в более реалистичных моделях добавляют внутривидовую
конкуренцию, насыщение хищника и другие эффекты.

\subsection{Эпидемиологическая модель SIR}

Пусть $S(t)$, $I(t)$ и $R(t)$ --- числа восприимчивых, инфицированных и
удалённых из эпидемического процесса особей, а
\[
N=S+I+R=\text{const}.
\]
Одна из стандартных форм модели:
\[
\boxed{
\begin{aligned}
\dot S&=-\beta\frac{SI}{N},\\
\dot I&=\beta\frac{SI}{N}-\gamma I,\\
\dot R&=\gamma I.
\end{aligned}
}
\]
Базовое репродуктивное число
\[
\boxed{R_0=\frac{\beta}{\gamma}}
\]
показывает среднее число вторичных заражений в полностью восприимчивой
популяции. В начале эпидемии
\[
\dot I>0
\quad\Longleftrightarrow\quad
R_0\frac{S_0}{N}>1.
\]
При $S_0\approx N$ это даёт известный порог $R_0>1$.

\subsection{Пространственные модели}

Если существен пространственный перенос, к кинетике добавляют диффузию.
Например, для вектора концентраций $u(x,t)$:
\[
\frac{\partial u}{\partial t}
=D\Delta u+F(u).
\]
Такие \textbf{реакционно-диффузионные системы} используются для описания
распространения биологических популяций, морфогенеза, пространственных
эпидемий и химико-биологических волн.

\subsection{Биоинформатика: выравнивание последовательностей}

Пусть даны две последовательности
\[
A=a_1\ldots a_n,
\qquad
B=b_1\ldots b_m.
\]
Задача глобального выравнивания состоит в поиске расположения совпадений,
замен и пропусков, максимизирующего суммарный балл. Она решается методом
динамического программирования. Для штрафа $d>0$ за пропуск и функции
сходства $s(a_i,b_j)$ рекуррентное соотношение Нидлмана--Вунша имеет вид
\[
\boxed{
F_{ij}=\max\left\{
\begin{array}{l}
F_{i-1,j-1}+s(a_i,b_j),\\
F_{i-1,j}-d,\\
F_{i,j-1}-d
\end{array}
\right\}.
}
\]
Граничные условия:
\[
F_{i0}=-id,
\qquad
F_{0j}=-jd.
\]
После заполнения таблицы оптимальное выравнивание восстанавливается
обратным проходом. Сложность классического алгоритма составляет
$O(nm)$ по времени. Для локального выравнивания алгоритм Смита--Уотермана
дополнительно допускает обнуление отрицательного значения.

\subsection{Другие модели биоинформатики}

В зависимости от задачи применяют вероятностные и статистические модели:
скрытые марковские модели для аннотации последовательностей, модели
эволюции для филогенетики, регрессионные и машинно-обучающие методы для
анализа высокоразмерных биологических данных. Важно, что математическая
модель всегда связана с конкретной биологической гипотезой и должна
проверяться на независимых данных.


\subsection{Точное решение логистической модели}

При $N(0)=N_0>0$ логистическое уравнение допускает явное решение
\[
    \boxed{
    N(t)=\frac{K}{1+\left(\dfrac{K}{N_0}-1\right)e^{-rt}}.
    }
\]
При $r>0$ оно монотонно стремится к $K$ при $0<N_0<K$ и убывает к $K$ при
$N_0>K$. Тем самым параметр $K$ определяет долгосрочный масштаб численности, а
$r$ --- характерное время выхода на насыщение.

\subsection{Интеграл модели Лотки--Вольтерры}

Для классической системы ``хищник--жертва'' существует первый интеграл
\[
    \boxed{
    H(x,y)=d x-c\ln x+b y-a\ln y=\text{const}.
    }
\]
Действительно, подстановка $\dot x=ax-bxy$, $\dot y=-cy+dxy$ даёт
$dH/dt=0$. В положительной четверти линии уровня $H$ образуют замкнутые кривые
вокруг равновесия
\[
    \left(\frac cd,\frac ab\right).
\]
Якобиан в этой точке имеет собственные значения
\[
    \lambda_{1,2}=\pm i\sqrt{ac},
\]
что согласуется с нейтральными колебаниями идеальной модели. Реальные системы
обычно требуют дополнительных механизмов затухания или насыщения.

\subsection{Порог и конечный размер эпидемии в SIR}

Обозначим базовое репродуктивное число через
\[
    \mathcal R_0=\frac{\beta}{\gamma}.
\]
Максимум числа инфицированных достигается, когда
\[
    \dot I=0
    \quad\Longleftrightarrow\quad
    \frac{S}{N}=\frac1{\mathcal R_0}.
\]
Кроме того,
\[
    \frac{dS}{dR}
    =-\frac{\beta}{\gamma N}S
    =-\frac{\mathcal R_0}{N}S,
\]
поэтому
\[
    \ln\frac{S(t)}{S(0)}
    =-\mathcal R_0\frac{R(t)-R(0)}{N}.
\]
При $t\to\infty$ это соотношение задаёт неявное уравнение для конечного размера
эпидемии. Оно показывает, что порог $\mathcal R_0>1$ определяет начальный рост,
но не означает заражение всей популяции.

\subsection{Реакционно--диффузионная устойчивость}

Для системы
\[
    u_t=D\Delta u+F(u)
\]
сначала исследуют однородное равновесие $F(u^*)=0$ по собственным значениям
якобиана $DF(u^*)$. Затем учитывают пространственные моды. Возможна ситуация,
когда равновесие устойчиво без диффузии, но некоторые пространственные моды
становятся неустойчивыми после её включения. Это механизм диффузионно
индуцированной, или тьюринговской, неустойчивости, применяемый в моделях
морфогенеза.

\subsection{Скрытые марковские модели в биоинформатике}

Скрытая марковская модель задаётся вероятностями переходов $a_{ij}$ между
скрытыми состояниями и вероятностями эмиссии $b_j(o)$. Если наблюдается
последовательность $o_1,\ldots,o_T$, прямой алгоритм использует рекурсию
\[
    \boxed{
    \alpha_t(j)
    =b_j(o_t)\sum_i\alpha_{t-1}(i)a_{ij}.
    }
\]
Вероятность наблюдаемой последовательности равна $\sum_j\alpha_T(j)$.
Алгоритм Витерби заменяет суммирование максимумом и находит наиболее вероятную
последовательность скрытых состояний. Такие модели естественно применяются к
аннотации последовательностей, поиску генов и распознаванию биологических
состояний вдоль цепочки ДНК или белка.

\subsection{Стохастические модели и ограниченность данных}

При малых численностях популяций непрерывная детерминированная модель может быть
недостаточна. Тогда используют процессы рождения--гибели или более общие
марковские модели, в которых число особей является случайной величиной. Выбор
между детерминированной и стохастической моделью определяется масштабом системы,
целевым вопросом и доступными данными. Параметры биологических моделей обычно
оцениваются по эксперименту, поэтому анализ чувствительности и проверка на
независимых данных являются частью самой модели, а не только постобработкой.

\subsection{Что важно сказать на экзамене}

Удобно привести три разных типа моделей: динамическую модель роста
(логистическое уравнение), систему взаимодействующих популяций
Лотки--Вольтерры и SIR. Для биоинформатики хорошим конкретным примером
является выравнивание последовательностей и рекуррентная формула
динамического программирования. Желательно уметь объяснить физический или
биологический смысл параметров и стационарных состояний.

\newpage
